\documentclass[12pt]{article}
\usepackage[T1]{fontenc}
\usepackage[utf8]{inputenc}
\usepackage{lmodern}
\usepackage{textcomp}
\usepackage{enumitem}
\usepackage{geometry}
\geometry{margin=1in}
\begin{document}
\begin{center}
Answer Key - Biology - Graduate Level Assessment 22 \
\small (Answers are ordered exactly as in the question paper.)
\end{center}

\section*{Multiple Choice}
\begin{enumerate}[label=\arabic*.]
\item \textbf{Answer:} A
\\ \textit{Options:} A) By assessing the fossil content found within the rock layers; B) By comparing the density of the rock to a known age database; C) By examining the color and texture of the rock; D) By counting the layers of the rock; E) By measuring the amount of carbon-14 remaining in the rock
\\ \textit{Note:} The correct answer is based on the principle that the fossil content within rock layers allows scientists to correlate and date the layers using the known ages of the fossils, thereby providing an estimate for the age of the rocks themselves.
\item \textbf{Answer:} A
\\ \textit{Options:} A) Animals use the sound of the ocean waves for directional cues in migration.; B) Animals rely solely on memory from previous migrations.; C) Animals migrate based on the patterns of the stars alone.; D) Animals follow changes in air pressure for migration.; E) Animals use temperature changes for migration.
\\ \textit{Note:} The correct answer highlights that many animals, particularly marine species, utilize auditory cues such as the sound of ocean waves to navigate during migration, demonstrating the importance of environmental sounds in their orientation and movement patterns.
\item \textbf{Answer:} A
\\ \textit{Options:} A) 108 nucleotides; B) 120 nucleotides; C) 150 nucleotides; D) 144 nucleotides; E) 60 nucleotides
\\ \textit{Note:} In a DNA molecule, the base pairing rules dictate that adenine pairs with thymine and cytosine pairs with guanine; since adenine constitutes 20 percent of the 180 base pairs, there are 36 adenine nucleotides, leading to 36 thymine nucleotides, which means the remaining 108 nucleotides must be cytosine and guanine, with an equal number of each, resulting in 54 cytosine nucleotides.
\item \textbf{Answer:} D
\\ \textit{Options:} A) RNase only; B) a DNA primer and RNase; C) a DNA or RNA primer and DNA polymerase; D) a DNA or RNA primer and reverse transcriptase; E) a RNA primer and DNA polymerase
\\ \textit{Note:} The synthesis of an RNA/DNA hybrid from a single-stranded RNA template requires a DNA or RNA primer to initiate the elongation process and reverse transcriptase to catalyze the synthesis of the complementary DNA strand.
\item \textbf{Answer:} A
\\ \textit{Options:} A) Antibiotic resistance arises from bacteria absorbing nutrients that counteract the drug; B) Resistance develops due to prolonged exposure to antibiotics; C) Bacteria develop a natural immunity to drugs over time; D) Bacteria develop resistance through spontaneous generation without exposure to antibiotics; E) Drug resistance is caused solely by the use of antibacterial hand sanitizers
\\ \textit{Note:} The correct answer highlights that bacteria can develop antibiotic resistance by taking up and utilizing specific nutrients or genetic material that enable them to neutralize or evade the effects of the antibiotic, demonstrating the adaptive nature of bacterial survival mechanisms.
\end{enumerate}

\section*{True / False}
\begin{enumerate}[label=\arabic*.]
\setcounter{enumi}{5}
\item \textbf{Answer:} False
\\ \textit{Options:} A) True; B) False
\\ \textit{Note:} In Southern Italy, hepatitis G virus infection is widespread among patients with chronic hepatitis, independently of parenteral risk factors. Its frequency in subjects with cryptogenic liver disease parallels that observed in hepatitis C virus chronic liver disease, thus ruling out an aetiologic role of hepatitis G virus. TT virus infection is uncommon in patients with cryptogenic or hepatitis C virus-related liver disease who do not have a history of parenteral exposure.
\item \textbf{Answer:} True
\\ \textit{Options:} A) True; B) False
\\ \textit{Note:} These findings indicate that a relationship between multiple sclerosis and streptococcal infections may exist, but to acquire a better understanding of the role of group A streptococci in the pathogenesis of multiple sclerosis, more studies with animal models are necessary.
\item \textbf{Answer:} True
\\ \textit{Options:} A) True; B) False
\\ \textit{Note:} Our study confirms that students are able to identify QA issues and propose solutions. We recommend that students have a formally recognized place in day surgery QA programmes, to close the QA loop and to adequately prepare them for medical practice in the 21 st century.
\item \textbf{Answer:} False
\\ \textit{Options:} A) True; B) False
\\ \textit{Note:} Baseline pain intensity does not predict the outcome after an appropriate opioid titration. It is likely that non-homogeneous pain treatment would have biased the outcome of a previous work.
\item \textbf{Answer:} True
\\ \textit{Options:} A) True; B) False
\\ \textit{Note:} These cases illustrate that tumours involving eyelids and parotid glands can present simultaneously or sequentially, and either of these structures could be the focus of primary or metastatic tumour. The important message for oculoplastic and parotid surgeons is to routinely assess both the periocular and parotid area when patients present with a mass in either structure.
\end{enumerate}

\section*{Long Form Response}
\begin{enumerate}[label=\arabic*.]
\setcounter{enumi}{10}
\item \textbf{Answer:} Low external pH plays a significant role in ATP synthesis by isolated mitochondria through its effect on the proton gradient across the inner mitochondrial membrane. A lower external pH increases the concentration of protons (H+) in the intermembrane space, enhancing the electrochemical gradient that drives the diffusion of H+ ions from the intermembrane space into the mitochondrial matrix. This increased diffusion rate facilitates the operation of ATP synthase, the enzyme responsible for synthesizing ATP, as it relies on the flow of protons to catalyze the conversion of ADP and inorganic phosphate into ATP. Consequently, a low external pH can boost ATP production efficiency by promoting a more favorable proton motive force, essential for effective oxidative phosphorylation.
\\ \textit{Note:} correct because a low external pH increases proton concentration in the intermembrane space, thereby enhancing the electrochemical gradient that drives H+ diffusion into the matrix, which is crucial for ATP synthase activity and ATP production.
\item \textbf{Answer:} Fermentation plays a crucial role in cellular respiration under anaerobic conditions by allowing cells to produce energy without oxygen. During this process, glucose is partially degraded, resulting in the production of ATP while also regenerating NAD+ from NADH. This regeneration of NAD+ is vital because it enables glycolysis to continue, thus facilitating the ongoing production of ATP in the absence of aerobic respiration. The significance of fermentation lies in its ability to sustain energy production in environments where oxygen is limited, thereby supporting cellular functions and survival. Ultimately, fermentation highlights the versatility of cellular metabolism in adapting to various environmental conditions.
\\ \textit{Note:} correct because it effectively describes how fermentation allows for energy production in the absence of oxygen by regenerating NAD+, which is essential for the continuation of glycolysis and ATP synthesis.
\end{enumerate}

\vfill
\noindent\textit{End of Answer Key}
\end{document}